\documentclass[12pt]{article}
\usepackage[utf8]{inputenc}
\usepackage[T1]{fontenc}
\usepackage[margin=1in]{geometry}
\usepackage{setspace}
\usepackage{xcolor}
\usepackage[colorlinks=true, urlcolor=blue!60!black, linkcolor=blue!60!black]{hyperref}
\setlength{\parindent}{2em}
\setlength{\parskip}{1.25em}
\renewcommand{\baselinestretch}{1.0}
\begin{document}
\begin{center}
{\LARGE \textit{Gibbons v. Ogden}, 22 U.S. (9 Wheat.) 1 (1824)\par}
\vspace{0.8em}
{\large Case Brief --- Four-Page Edition\par}
\end{center}
\vspace{1em}

\section{Memory Jogger}

Congress's Commerce Clause power encompasses interstate navigation and is plenary within its sphere, so a federal coasting license preempts New York's conflicting steamboat monopoly under the Supremacy Clause (\href{https://www.courtlistener.com/opinion/85412/gibbons-v-ogden/}{22 U.S.\ 1}).

\section{Detailed Case Facts}

Through statutes enacted from 1798 to 1811, New York granted Robert R.\ Livingston and Robert Fulton the exclusive right to navigate all New York waters with steam-propelled vessels for up to thirty years, enforced by forfeiture of unlicensed steamboats. The right to run steamboats between Elizabethtown, New Jersey, and New York City passed by assignment to Aaron Ogden, a former Governor of New Jersey.

Thomas Gibbons, Ogden's former business partner, began running two rival steamboats, the \textit{Stoudinger} and the \textit{Bellona}, on the same route. He had no license from the monopoly, but his vessels were enrolled and licensed for the coasting trade under the federal Coasting Act of 1793. Neighboring states had retaliated against the monopoly: Connecticut barred New York-licensed vessels, and New Jersey passed an ``act of retortion'' authorizing damages against anyone enforcing it. In 1818 Ogden sued in New York's Court of Chancery to enjoin Gibbons.

\section{Procedural History}

Chancellor James Kent enjoined Gibbons, holding the monopoly a valid exercise of concurrent state power and the federal license a mere grant of national character. New York's Court for the Trial of Impeachments and Correction of Errors unanimously affirmed. A first appeal was dismissed for want of a final decree. After final decree, Gibbons appealed again; with Daniel Webster and Attorney General William Wirt arguing for him, the Court on March 2, 1824, reversed and ordered Ogden's bill dismissed.

\section{Judicial Votes}

Reversed, \textbf{6--0}.
\begin{itemize}
\item \textbf{Majority:} Chief Justice Marshall, joined by Justices Washington, Todd, Duvall, and Story.
\item \textbf{Concurring in the judgment:} Justice William Johnson.
\item \textbf{Not participating:} Justice Smith Thompson.
\end{itemize}

\section{Holding}

``Commerce'' includes every species of commercial intercourse, including navigation by steamboats carrying passengers. Commerce ``among the several States'' does not stop at state lines but does not reach commerce ``completely internal'' to one state. Congress's power is plenary, subject only to constitutional limits, and a federal coasting license affirmatively authorizes interstate navigation, preempting the conflicting monopoly under the Supremacy Clause. Marshall's sympathy for Webster's exclusivity (dormant Commerce Clause) theory, his description of state police-power legislation, and the patent-clause question are dicta.

\section{Analysis of Opinions}

\textbf{Majority (Marshall).} Marshall rejected ``strict construction'' that would ``cripple the government,'' giving the words their natural meaning in light of the Constitution's objects. ``Commerce'' is not mere buying and selling but ``commercial intercourse'' in all its branches, including navigation, as the Port Preference Clause and the embargo and navigation acts confirm. ``Among'' means ``intermingled with,'' so federal power reaches into a state's interior whenever commerce concerns more than one state, while the enumeration presupposes the ``exclusively internal commerce of a State.'' ``To regulate'' means to prescribe the rule, and the power's only restraints, beyond express prohibitions, are political checks through elections. State inspection, quarantine, and health laws flow from the police power, not a concurrent commerce power. Rather than decide whether the Commerce Clause alone voided the monopoly, Marshall turned to the Supremacy Clause: the Coasting Act license grants permission to carry on the coasting trade, with cargo or passengers, by wind or steam, so the state statutes were void to the extent of the collision.

\textbf{Johnson (concurring in the judgment).} Johnson would rest on the Constitution itself. Drawing on the law of nations and the purpose of ending interstate trade wars, he argued that the commerce power passed entirely to Congress and is exclusive, leaving states no authority over interstate navigation even absent the Coasting Act.

\section{Examples: Future Applications}

\begin{enumerate}
\item \textbf{Same-Side: Mississippi River Ferry Monopoly.} An Illinois exclusive franchise for electric vessels cannot be enforced against a federally licensed Missouri operator; the license governs regardless of propulsion technology.
\item \textbf{Same-Side: Autonomous Trucks.} A state's exclusion of federally licensed autonomous freight trucks to protect local firms collides with valid federal regulation and must yield.
\item \textbf{Opposite-Side: Livestock Inspection.} A non-discriminatory health inspection law for local livestock, with no conflicting federal statute, falls within the states' ``immense mass of legislation.''
\item \textbf{Opposite-Side: Rowboats on a Landlocked Lake.} Licensing recreational rowboats on a lake within one county regulates commerce ``completely internal'' to the state.
\item \textbf{Fence-Sitter: Intrastate Hudson Ferry.} An exclusive franchise for an Albany--New York City ferry carrying many out-of-state travelers, with no governing federal statute, is uncertain: Marshall reserved internal commerce and relied on statutory collision, but Johnson's exclusivity theory points toward invalidation.
\end{enumerate}

\section{Critique}

\textbf{Originalist views.} Randy Barnett argues that ``commerce'' originally meant trade and transportation of goods, including navigation, supporting \textit{Gibbons}'s result but not later extensions to manufacturing or agriculture. Jack Balkin reads ``intercourse'' as broader social and economic interaction, a reading Barnett and Heritage commentators contest. Textualists also question Marshall's gloss on ``among'' and his dicta favoring an unenumerated exclusive power.

\textbf{Progressive and doctrinal views.} Many celebrate \textit{Gibbons} as the cornerstone of the national common market and modern regulatory authority. Norman Williams criticizes Marshall's ``hedge'' of retreating to a strained reading of the Coasting Act, which left decades of confusion until \textit{Cooley v.\ Board of Wardens} (1851). Marshall's twin commitments, a power ``complete in itself'' and a reserved internal commerce, resurfaced in \textit{Wickard}, \textit{Lopez}, and \textit{Raich}.

\section{Key Quotations}

\begin{quote}
``It is the power to regulate; that is, to prescribe the rule by which commerce is to be governed. This power, like all others vested in Congress, is complete in itself, may be exercised to its utmost extent, and acknowledges no limitations, other than are prescribed in the Constitution.''
\end{quote}
\noindent --- Chief Justice Marshall, 22 U.S.\ at 196.

\begin{quote}
``Commerce, undoubtedly, is traffic, but it is something more: it is intercourse. It describes the commercial intercourse between nations, and parts of nations, in all its branches, and is regulated by prescribing rules for carrying on that intercourse.''
\end{quote}
\noindent --- Chief Justice Marshall, 22 U.S.\ at 189--90.

\begin{quote}
``The enumeration presupposes something not enumerated; and that something, if we regard the language or the subject of the sentence, must be the exclusively internal commerce of a State.''
\end{quote}
\noindent --- Chief Justice Marshall, 22 U.S.\ at 195.

\end{document}
